\begin{proposition}[Physical invariance under ambiguity-triviality]
\label{prop:Main-physical-invariance-ambiguity}
Fix an accepted tick \(k\), and let \(J_1,J_2\in[\Xi^{(1)}_{\Lambda_k}]\), so
\(J_2-J_1\in\mathcal A^-_{\Lambda_k}\). Define
\[
\delta^{(i)}_{\Lambda_k}:=\{J_i,-\}_{\mathrm P,\Lambda_k}
\qquad (i=1,2).
\]
If every ambiguity direction \(A\in\mathcal A^-_{\Lambda_k}\) acts trivially on the
physical observable sector, then \(\delta^{(1)}_{\Lambda_k}\) and
\(\delta^{(2)}_{\Lambda_k}\) agree there. More generally, if for every
\(A\in\mathcal A^-_{\Lambda_k}\) and every physical local observable \(F\) the
Peierls bracket \(\{A,F\}_{\mathrm P,\Lambda_k}\) vanishes in the physical quotient
\textnormal{(equivalently: belongs to the physical null ideal)}, then the induced
physical Peierls derivation is independent of the chosen representative in
\([\Xi^{(1)}_{\Lambda_k}]\).
\end{proposition}

\begin{proof}
Since \(J_2-J_1\in\mathcal A^-_{\Lambda_k}\), bilinearity of the Peierls bracket gives
\[
\delta^{(2)}_{\Lambda_k}(F)-\delta^{(1)}_{\Lambda_k}(F)
=\{J_2-J_1,F\}_{\mathrm P,\Lambda_k}.
\]
If every ambiguity direction acts trivially on the physical observable sector, the
right-hand side vanishes there, so the two derivations agree on physical observables.
Under the weaker null-ideal hypothesis, the same identity shows that the difference
vanishes after passage to the physical quotient, hence the induced physical derivation
depends only on the class \([\Xi^{(1)}_{\Lambda_k}]\).
\end{proof}